\begin{lemma}
\label{p2:lem:ordinary-tangent-correction}
There is an exact sequence
\begin{equation}
 0\longrightarrow T_X(-\log W)
 \longrightarrow T_X
 \longrightarrow\nu_*N_\nu
 \longrightarrow0,
 \label{p2:eq:ordinary-tangent-sequence}
\end{equation}
and
\begin{equation}
 N_\nu\simeq K_{E_0}.
 \label{p2:eq:cusp-normal-canonical}
\end{equation}
\end{lemma}

\begin{proof}
The last arrow in \eqref{p2:eq:ordinary-tangent-sequence} is the canonical
composite
\[
 T_X\longrightarrow\nu_*\nu^*T_X
 \longrightarrow\nu_*N_\nu.
\]
At a point where $r\in\{1,2,3\}$ branches of
$W=(z_0\cdots z_{r-1}=0)$ meet, it is
\[
 \sum_i a_i\partial_{z_i}
 \longmapsto
 \bigl(a_i|_{z_i=0}\bigr)_{0\le i<r}.
\]
The normal coefficients can be chosen independently, so this map is
surjective.  Its kernel consists of the vector fields preserving every
branch ideal $(z_i)$, namely $T_X(-\log W)$.  This local description also
shows why the quotient is the pushforward of the branchwise normal line on
the normalization, rather than the single sheaf $\cO_W(W)$.

On the dense cusp torus consider
\begin{equation}
 \Theta=
 \frac{dx_1}{x_1}\wedge\frac{dx_2}{x_2}\wedge dq.
 \label{p2:eq:cusp-three-form}
\end{equation}
In a maximal toric chart the three ray generators form a lattice basis and
$q=z_0z_1z_2$.  The associated character matrix has determinant $\pm1$,
so $\Theta=\pm dz_0\wedge dz_1\wedge dz_2$.  Thus $\Theta$ extends
holomorphically and without zeros across every toric boundary stratum.
Under the full cusp action,
$d\log x_i$ changes only by a multiple of $dq$ or $d\log q$; wedging with
$dq$ removes every added term.  Hence $\Theta$ is invariant under all
$\lambda\in\overline\Lambda$ and descends to a nowhere-zero section of
$K_X$ on a neighbourhood of $W$.

The adjunction formula for the immersion $\nu:E_0\to X$ gives
\[
 K_{E_0}\simeq\nu^*K_X\otimes N_\nu.
\]
The descended form trivializes $\nu^*K_X$, proving
\eqref{p2:eq:cusp-normal-canonical}.
\end{proof}